\section{AHSME 1971}
        \begin{enumerate}
        	\item (\textit{AHSME 1971}) The number of digits in the number $N=2^{12}\times 5^8$ is


 $\textbf{(A) }9\qquad \textbf{(B) }10\qquad \textbf{(C) }11\qquad \textbf{(D) }12\qquad  \textbf{(E) }20$


	\item (\textit{AHSME 1971}) If $b$ men take $c$ days to lay $f$ bricks, then the number of days it will take $c$ men working at the same rate to lay $b$ bricks, is


 $\textbf{(A) }fb^2\qquad \textbf{(B) }b/f^2\qquad \textbf{(C) }f^2/b\qquad \textbf{(D) }b^2/f\qquad  \textbf{(E) }f/b^2$


	\item (\textit{AHSME 1971}) If the point $(x,-4)$ lies on the straight line joining the points $(0,8)$ and $(-4,0)$ in the $xy$-plane, then $x$ is equal to


 $\textbf{(A) }-2\qquad \textbf{(B) }2\qquad \textbf{(C) }-8\qquad \textbf{(D) }6\qquad  \textbf{(E) }-6$


	\item (\textit{AHSME 1971}) After simple interest for two months at $5$\% per annum was credited, a Boy Scout Troop had a total 
of $\textdollar 255.31$ in the Council Treasury. The interest credited was a number of dollars plus the following number of cents


 $\textbf{(A) }11\qquad \textbf{(B) }12\qquad \textbf{(C) }13\qquad \textbf{(D) }21\qquad  \textbf{(E) }31$


	\item (\textit{AHSME 1971}) Points $A,B,Q,D$, and $C$ lie on the circle shown and the measures of arcs $\widehat{BQ}$ and $\widehat{QD}$ 
are $42^\circ$ and $38^\circ$ respectively. The sum of the measures of angles $P$ and $Q$ is


 $\textbf{(A) }80^\circ\qquad \textbf{(B) }62^\circ\qquad \textbf{(C) }40^\circ\qquad \textbf{(D) }46^\circ\qquad  \textbf{(E) }\text{None of these}$

                    \begin{figure}[H]
                        \centering
                        \includegraphics[width=0.5\textwidth]{Problem/AHSME_I_1971_5_Problem_Figure_1.png}
                    \end{figure}
                    


	\item (\textit{AHSME 1971}) Let $\ast$ be the symbol denoting the binary operation on the set $S$ of all non-zero real numbers as follows: 
For any two numbers $a$ and $b$ of $S$, $a\ast b=2ab$. Then the one of the following statements which is not true, is


 $\textbf{(A) }\ast\text{ is commutative over }S \qquad \textbf{(B) }\ast\text{ is associative over }S\qquad \\ \textbf{(C) }\frac{1}{2}\text{ is an identity element for }\ast\text{ in }S\qquad  \textbf{(D) }\text{Every element of }S\text{ has an inverse for }\ast\qquad  \textbf{(E) }\dfrac{1}{2a}\text{ is an inverse for }\ast\text{ of the element }a\text{ of }S$


	\item (\textit{AHSME 1971}) $2^{-(2k+1)}-2^{-(2k-1)}+2^{-2k}$ is equal to 


 $\textbf{(A) }2^{-2k}\qquad \textbf{(B) }2^{-(2k-1)}\qquad \textbf{(C) }-2^{-(2k+1)}\qquad \textbf{(D) }0\qquad  \textbf{(E) }2$


	\item (\textit{AHSME 1971}) The solution set of $6x^2+5x<4$ is the set of all values of $x$ such that


 $\textbf{(A) }\textstyle -2<x<1\qquad \textbf{(B) }-\frac{4}{3}<x<\frac{1}{2}\qquad \textbf{(C) }-\frac{1}{2}<x<\frac{4}{3}\qquad \\ \textbf{(D) }x<\textstyle\frac{1}{2}\text{ or }x>-\frac{4}{3}\qquad \textbf{(E) }x<-\frac{4}{3}\text{ or }x>\frac{1}{2}$


	\item (\textit{AHSME 1971}) An uncrossed belt is fitted without slack around two circular pulleys with radii of $14$ inches and $4$ inches. 
If the distance between the points of contact of the belt with the pulleys is $24$ inches, then the distance 
between the centers of the pulleys in inches is


 $\textbf{(A) }24\qquad \textbf{(B) }2\sqrt{119}\qquad \textbf{(C) }25\qquad \textbf{(D) }26\qquad  \textbf{(E) }4\sqrt{35}$


	\item (\textit{AHSME 1971}) Each of a group of $50$ girls is blonde or brunette and is blue eyed or brown eyed. If $14$ are blue-eyed blondes, 
$31$ are brunettes, and $18$ are brown-eyed, then the number of brown-eyed brunettes is


 $\textbf{(A) }5\qquad \textbf{(B) }7\qquad \textbf{(C) }9\qquad \textbf{(D) }11\qquad  \textbf{(E) }13$


	\item (\textit{AHSME 1971}) The numeral $47$ in base a represents the same number as $74$ in base $b$. Assuming that both bases are positive 
integers, the least possible value of $a+b$ written as a Roman numeral, is


 $\textbf{(A) }\mathrm{XIII}\qquad \textbf{(B) }\mathrm{XV}\qquad \textbf{(C) }\mathrm{XXI}\qquad \textbf{(D) }\mathrm{XXIV}\qquad \textbf{(E) }\mathrm{XVI}$


	\item (\textit{AHSME 1971}) For each integer $N>1$, there is a mathematical system in which two or more positive integers are defined 
to be congruent if they leave the same non-negative remainder when divided by $N.$ If $69,90$, and $125$ are 
congruent in one such system, then in that same system, $81$ is congruent to


 $\textbf{(A) }3\qquad \textbf{(B) }4\qquad \textbf{(C) }5\qquad \textbf{(D) }7\qquad  \textbf{(E) }8$


	\item (\textit{AHSME 1971}) If $(1.0025)^{10}$ is evaluated correct to $5$ decimal places, then the digit in the fifth decimal place is


 $\textbf{(A) }0\qquad \textbf{(B) }1\qquad \textbf{(C) }2\qquad \textbf{(D) }5\qquad  \textbf{(E) }8$


	\item (\textit{AHSME 1971}) The number $(2^{48}-1)$ is exactly divisible by two numbers between $60$ and $70$. These numbers are


 $\textbf{(A) }61,63\qquad \textbf{(B) }61,65\qquad \textbf{(C) }63,65\qquad \textbf{(D) }63,67\qquad  \textbf{(E) }67,69$


	\item (\textit{AHSME 1971}) An aquarium on a level table has rectangular faces and is $10$ inches wide and $8$ inches high. When it was tilted, 
the water in it covered an $8\times 10$ end but only $\frac{3}{4}$ of the rectangular bottom. The depth of the 
water when the bottom was again made level, was


 $\textbf{(A) }2\textstyle{\frac{1}{2}}"\qquad \textbf{(B) }3"\qquad \textbf{(C) }3\textstyle{\frac{1}{4}}"\qquad \textbf{(D) }3\textstyle{\frac{1}{2}}"\qquad \textbf{(E) }4"$


	\item (\textit{AHSME 1971}) After finding the average of $35$ scores, a student carelessly included the average with the $35$ scores and found the 
average of these $36$ numbers. The ratio of the second average to the true average was


 $\textbf{(A) }1:1\qquad \textbf{(B) }35:36\qquad \textbf{(C) }36:35\qquad \textbf{(D) }2:1\qquad  \textbf{(E) }\text{None of these}$


	\item (\textit{AHSME 1971}) A circular disk is divided by $2n$ equally spaced radii($n>0$) and one secant line. The maximum number of non-overlapping 
areas into which the disk can be divided is


 $\textbf{(A) }2n+1\qquad \textbf{(B) }2n+2\qquad \textbf{(C) }3n-1\qquad \textbf{(D) }3n\qquad  \textbf{(E) }3n+1$


	\item (\textit{AHSME 1971}) The current in a river is flowing steadily at $3$ miles per hour. A motor boat which travels at a constant rate in still 
water goes downstream $4$ miles and then returns to its starting point. The trip takes one hour, excluding the time spent 
in turning the boat around. The ratio of the downstream to the upstream rate is 


 $\textbf{(A) }4:3\qquad \textbf{(B) }3:2\qquad \textbf{(C) }5:3\qquad \textbf{(D) }2:1\qquad  \textbf{(E) }5:2$


	\item (\textit{AHSME 1971}) If the line $y=mx+1$ intersects the ellipse $x^2+4y^2=1$ exactly once, then the value of $m^2$ is


 $\textbf{(A) }\textstyle\frac{1}{2}\qquad \textbf{(B) }\frac{2}{3}\qquad \textbf{(C) }\frac{3}{4}\qquad \textbf{(D) }\frac{4}{5}\qquad \textbf{(E) }\frac{5}{6}$


	\item (\textit{AHSME 1971}) The sum of the squares of the roots of the equation $x^2+2hx=3$ is $10$. The absolute value of $h$ is equal to


 $\textbf{(A) }-1\qquad \textbf{(B) }\textstyle\frac{1}{2}\qquad \textbf{(C) }\textstyle\frac{3}{2}\qquad \textbf{(D) }2\qquad \textbf{(E) }\text{None of these}$


	\item (\textit{AHSME 1971}) If $\log_2(\log_3(\log_4 x))=\log_3(\log_4(\log_2 y))=\log_4(\log_2(\log_3 z))=0$, then the sum $x+y+z$ is equal to


 $\textbf{(A) }50\qquad \textbf{(B) }58\qquad \textbf{(C) }89\qquad \textbf{(D) }111\qquad  \textbf{(E) }1296$


	\item (\textit{AHSME 1971}) If $w$ is one of the imaginary roots of the equation $x^3=1$, then the product $(1-w+w^2)(1+w-w^2)$ is equal to


 $\textbf{(A) }4\qquad \textbf{(B) }w\qquad \textbf{(C) }2\qquad \textbf{(D) }w^2\qquad  \textbf{(E) }1$


	\item (\textit{AHSME 1971}) Teams $A$ and $B$ are playing a series of games. If the odds for either to win any game are even and 
Team $A$ must win two or Team $B$ three games to win the series, then the odds favoring Team $A$ to win the series are


 $\textbf{(A) }11\text{ to }5\qquad \textbf{(B) }5\text{ to }2\qquad \textbf{(C) }8\text{ to }3\qquad \textbf{(D) }3\text{ to }2\qquad \textbf{(E) }13\text{ to }6$


	\item (\textit{AHSME 1971}) 
                    \begin{figure}[H]
                        \centering
                        \includegraphics[width=0.5\textwidth]{Problem/AHSME_I_1971_24_Problem_Figure_1.png}
                    \end{figure}
                    


 Pascal's triangle is an array of positive integers(See figure), in which the first row is $1$, the second row is two $1$'s, 
each row begins and ends with $1$, and the $k^\text{th}$ number in any row when it is not $1$, is the sum of the $k^\text{th}$
and $(k-1)^\text{th}$ numbers in the immediately preceding row. The quotient of the number of numbers in the first $n$ rows 
which are not $1$'s and the number of $1$'s is


 $\textbf{(A) }\dfrac{n^2-n}{2n-1}\qquad \textbf{(B) }\dfrac{n^2-n}{4n-2}\qquad \textbf{(C) }\dfrac{n^2-2n}{2n-1}\qquad \textbf{(D) }\frac{n^2-3n+2}{4n-2}\qquad \textbf{(E) }\text{None of these}$


	\item (\textit{AHSME 1971}) A teen age boy wrote his own age after his father's. From this new four place number, he subtracted the absolute value of the 
difference of their ages to get $4,289$. The sum of their ages was


 $\textbf{(A) }48\qquad \textbf{(B) }52\qquad \textbf{(C) }56\qquad \textbf{(D) }59\qquad  \textbf{(E) }64$


	\item (\textit{AHSME 1971}) 
                    \begin{figure}[H]
                        \centering
                        \includegraphics[width=0.5\textwidth]{Problem/AHSME_I_1971_26_Problem_Figure_1.png}
                    \end{figure}
                    


 In $\triangle ABC$, point $F$ divides side $AC$ in the ratio $1:2$. Let $E$ be the point of intersection of side $BC$ and $AG$ where $G$ is the 
midpoint of $BF$. The point $E$ divides side $BC$ in the ratio


 $\textbf{(A) }1:4\qquad \textbf{(B) }1:3\qquad \textbf{(C) }2:5\qquad \textbf{(D) }4:11\qquad  \textbf{(E) }3:8$


	\item (\textit{AHSME 1971}) A box contains chips, each of which is red, white, or blue. The number of blue chips is at least half the number of white chips, 
and at most one third the number of red chips. The number which are white or blue is at least $55$. The minimum number of red chips is


 $\textbf{(A) }24\qquad \textbf{(B) }33\qquad \textbf{(C) }45\qquad \textbf{(D) }54\qquad  \textbf{(E) }57$


	\item (\textit{AHSME 1971}) Nine lines parallel to the base of a triangle divide the other sides each into $10$ equal segments and the area into $10$ distinct parts. 
If the area of the largest of these parts is $38$, then the area of the original triangle is 


 $\textbf{(A) }180\qquad \textbf{(B) }190\qquad \textbf{(C) }200\qquad \textbf{(D) }210\qquad  \textbf{(E) }240$


	\item (\textit{AHSME 1971}) Given the progression $10^{\dfrac{1}{11}}, 10^{\dfrac{2}{11}}, 10^{\dfrac{3}{11}}, 10^{\dfrac{4}{11}},\dots , 10^{\dfrac{n}{11}}$. 
The least positive integer $n$ such that the product of the first $n$ terms of the progression exceeds $100,000$ is


 $\textbf{(A) }7\qquad \textbf{(B) }8\qquad \textbf{(C) }9\qquad \textbf{(D) }10\qquad  \textbf{(E) }11$


	\item (\textit{AHSME 1971}) Given the linear fractional transformation of $x$ into $f_1(x)=\dfrac{2x-1}{x+1}$. Define $f_{n+1}(x)=f_1(f_n(x))$ for $n=1,2,3,\cdots$. 
Assuming that $f_{35}(x)=f_5(x)$, it follows that $f_{28}(x)$ is equal to


 $\textbf{(A) }x\qquad \textbf{(B) }\frac{1}{x}\qquad \textbf{(C) }\frac{x-1}{x}\qquad \textbf{(D) }\frac{1}{1-x}\qquad  \textbf{(E) }\text{None of these}$


	\item (\textit{AHSME 1971}) 
                    \begin{figure}[H]
                        \centering
                        \includegraphics[width=0.5\textwidth]{Problem/AHSME_I_1971_31_Problem_Figure_1.png}
                    \end{figure}
                    


 Quadrilateral $ABCD$ is inscribed in a circle with side $AD$, a diameter of length $4$. If sides $AB$ and $BC$ each have length $1$, 
then side $CD$ has length


 $\textbf{(A) }\frac{7}{2}\qquad \textbf{(B) }\frac{5\sqrt{2}}{2}\qquad \textbf{(C) }\sqrt{11}\qquad \textbf{(D) }\sqrt{13}\qquad \textbf{(E) }2\sqrt{3}$


	\item (\textit{AHSME 1971}) If $s=(1+2^{-\frac{1}{32}})(1+2^{-\frac{1}{16}})(1+2^{-\frac{1}{8}})(1+2^{-\frac{1}{4}})(1+2^{-\frac{1}{2}})$, then $s$ is equal to


 $\textbf{(A) }\textstyle{\frac{1}{2}}(1-2^{-\frac{1}{32}})^{-1}\qquad \textbf{(B) }(1-2^{-\frac{1}{32}})^{-1}\qquad \textbf{(C) }1-2^{-\frac{1}{32}}\qquad \\ \textbf{(D) }\textstyle{\frac{1}{2}}(1-2^{-\frac{1}{32}})\qquad  \textbf{(E) }\frac{1}{2}$


	\item (\textit{AHSME 1971}) If $P$ is the product of $n$ quantities in Geometric Progression, $S$ their sum, and $S'$ the sum of their reciprocals, 
then $P$ in terms of $S, S'$, and $n$ is


 $\textbf{(A) }(SS')^{\frac{1}{2}n}\qquad \textbf{(B) }(S/S')^{\frac{1}{2}n}\qquad \textbf{(C) }(SS')^{n-2}\qquad \textbf{(D) }(S/S')^n\qquad \textbf{(E) }(S/S')^{\frac{1}{2}(n-1)}$


	\item (\textit{AHSME 1971}) An ordinary clock in a factory is running slow so that the minute hand passes the hour hand at the usual dial position($12$ o'clock, etc.) 
but only every $69$ minutes. At time and one-half for overtime, the extra pay to which a $\textdollar 4.00$ per hour worker should be entitled 
after working a normal $8$ hour day by that slow running clock, is


 $\textbf{(A) }\textdollar 2.30\qquad \textbf{(B) }\textdollar 2.60\qquad \textbf{(C) }\textdollar 2.80\qquad \textbf{(D) }\textdollar 3.00\qquad \textbf{(E) }\textdollar 3.30$


	\item (\textit{AHSME 1971}) Each circle in an infinite sequence with decreasing radii is tangent externally to the one following it and to both sides of a given right angle. 
The ratio of the area of the first circle to the sum of areas of all other circles in the sequence, is


 $\textbf{(A) }(4+3\sqrt{2}):4\qquad \textbf{(B) }9\sqrt{2}:2\qquad \textbf{(C) }(16+12\sqrt{2}):1\qquad \\ \textbf{(D) }(2+2\sqrt{2}):1\qquad  \textbf{(E) }(3+2\sqrt{2}):1$


\end{enumerate}